\section{Discussion and Limitations}\label{sec:disc}

\noindent\textbf{Why this coupling did not pay here.}
The path-aware cost replaces $d(p^r,p_\tau)$ by the arrival time
$A^{\mathcal{B}}(p^r,p_\tau,t_0)$. Within the window the two differ only by
the waits and detours that reservations impose on the empty trip to the
pickup; beyond the window \cref{alg:astar} appends a shortest path that
ignores reservations. In our
layouts that difference is small on two counts. First, the empty trip is a
minor part of a task: under batch load with 32 robots, service time is
282--287 ticks of which 216--220 are queueing before assignment
(\cref{tab:main}); the remaining 66 ticks or so contain the trip to the
pickup, two dwells and the loaded trip, and the pickup leg is the shorter one because
robots start near the aisles. Second, with 2-wide aisles the planner absorbs
most interactions by a one-tick wait or a lane change: a run of 200 tasks
executes 274 wait actions against 14.9 thousand moves, about 1.4 waits per
task. We did not record how often the two costs lead to different
matchings, so the next step is an inference, not a measurement: the
arrival-time estimate should rarely re-order candidates, and when it does the
change is within the noise of task sampling. Congestion that
does matter concentrates at the docks, which the pickup-leg cost cannot see
and which the endpoint capacity already regulates: on the warehouse layout,
the nine docks of capacity three bound the number of tasks in flight
(\cref{subsec:main}), and, as noted above, most of the service time under
batch load is spent waiting before assignment. This is in line with
\citet{kou2020idle}, who assign robots to stations in sortation centres so
as to minimise station idle time.
We expected one-lane
aisles, hotspot demand, long empty trips or saturated chargers to change the
balance and tested them (\cref{subsec:coupling}); the three inside the safe
envelope did not, within intervals of one to two percent, and the fourth is
confounded by stalled runs. The explanation above survives the test only in
a narrower form: even where arrival times deviate from free-space distances,
the assignments that minimise either one deliver throughputs that we cannot
tell apart, which is consistent with the queue in front of the docks, not
the trip to the pickup, setting the pace. The negative result remains a statement about our grid warehouses,
fleets of at most 48 robots, cost models and endpoint rules, not about
coupling in general. It is compatible with the positive results of prior
and concurrent work, which concern other signals or regimes: RMCA and
Detour-Delta price complete trips under planned paths or committed
reservations \citep{chen2021integrated,hirayama2026detour}; LNS-wPBS, whose
Hungarian-based insertion beats a greedy heuristic, assigns multi-goal task
sequences on estimates of service time \citep{xu2022multigoal}, whereas
\sname\ assigns one task per idle robot under a dock-capacity gate, so the
two greedy-versus-Hungarian comparisons are made on different problems; the
evaluations of flow-based assignment and GRAND include much larger fleets
\citep{zhang2025flow,gaber2026grand}; and the assigner of
\citet{zhu2026combined} chooses a robot's delivery port rather than the
robot for a given task. A coupling that also prices the loaded leg or the
dock queue, or that plans the assigned robots jointly as CBS-TA does for
one-shot instances \citep{honig2018cbsta}, is not ruled out.

\noindent\textbf{Cost analysis.}
The whole stack runs at 0.02--0.97 ms of planner time per simulated tick for 4
to 48 robots on one core (\cref{fig:scaling}); a tick of one second of robot
motion therefore leaves three orders of magnitude of headroom, and the
coupling's extra A* calls ($+35\%$ at 48 robots) are affordable even though
they are not useful here. Windowed CBS at the small fleet sizes where it is
exact costs 0.09--0.34 ms/tick. Memory is dominated by the cached distance
tables, one $|V|$-vector per endpoint.

\noindent\textbf{What made the loop robust.}
During development, each of three mechanisms removed a jam we had observed;
the second round of experiments measures all three (\cref{subsec:cascade},
\cref{subsec:ablation}). (i) The hold
cascade (\cref{alg:pp}) is how \sname\ makes ``no path found'' safe: without
some repair of the robots whose reservations cross a held cell, a hold
silently invalidates those reservations. The ablation shows that the
conflicts come from leaving invalidated reservations unrepaired: no repair
executes conflicts in every setting, and priority restarts, which fall back
to no repair after three attempts, do so in three of the four settings
(\cref{tab:cascade}); it does not compare re-planning with the IStay or
IAvoid remedy of \citet{morag2023adapting}. (ii) Endpoint capacities bound the queue in front of a
dock; with unbounded queues, robots waiting for the same dock filled both
lanes of the corridor and the robot leaving the dock could not get out.
(iii) Planning held robots first, across planning groups, resolves the one
gridlock that survived (i) and (ii): a robot that had just delivered at a
corner dock parks on the dock's single access cell, is boxed in by the robots
queuing for that dock, fails to plan, and---because parking robots are
planned last---never gets the priority it needs to leave. Without the
cascade, executed conflicts appear in most runs; with static instead of
dynamic priorities, the fleet gridlocks in 113 of 120 runs
(\cref{tab:cascade}); tightening (ii) from 3 to 2 robots per dock costs a
fifth of the throughput, and loosening it to 4 gains 9\% at twice the
holds. Failure handling in lifelong planning has been studied explicitly
\citep{morag2023adapting,yan2026lsmart,tang2026energy}, and (i) is a
variant of the IStay policy of \citet{morag2023adapting}; we report the
three mechanisms because each one was found by a jam in our own runs before
it was fixed, and \cref{subsec:cascade} measures them for our planner.

\noindent\textbf{Limitations.}
Everything in this paper is simulation on a unit-speed 4-connected grid with
synchronous ticks: no kinematics, no localisation error, no communication
delay, no real robots, so our grid-level results may not carry over to real
execution; \citet{varambally2022which} give a controlled study of how MAPF
models of different fidelity affect warehouse throughput under more realistic
execution. Tasks are sampled uniformly over endpoints except in
two coupling regimes; real demand is also time-varying, which we do not
model. The battery model is linear and the
charging policy is a fixed hysteresis; we do not optimise charging
schedules; under doubled drain it strands robots whenever the slots
saturate (\cref{tab:stress}), and \cref{prop:charge} only says which delay
bounds would prevent that, not how to enforce them. A charging policy that
reserves slots ahead of time or raises $\theta_{\mathrm{lo}}$ with the
charger queue is the obvious next step. Five seeds per configuration in the
first round bound its resolution to differences of a few percent; the second
round uses 30 to 60 seeds, and neither round was designed as an equivalence
test. The coupling study has no baseline from the coupling literature: we do
not run RMCA, LNS-wPBS or Detour-Delta, nor a cost that prices the whole
pickup-to-delivery trip, and we have no congested setting without the
dock-capacity gate or with fleets of hundreds of robots. These runs are the
evaluation we plan next; until then the null result says nothing about those
methods. We also did not record how often the static and path-aware costs
produce different matchings, which would test the explanation above
directly. Basic CBS, with PP for robots that share a goal, runs on groups of up to 32
robots with few node-budget fall-backs but gains only 1--3\% makespan there; symmetry reasoning
\citep{li2019symmetry} or bounded-suboptimal search \citep{li2021eecbs}
might gain more. Prioritized planning with a hold cascade is safe but not
complete: a robot can be delayed for several epochs, and although every run
with normal batteries finished, we have no proof that none can deadlock or
livelock. We did not run an IStay- or IAvoid-style variant, in which the
affected robots stay in place (under IAvoid, after at most one avoiding
move), so the ablation does not show that re-planning them is better than
making them stay. Our token-passing baselines follow \citet{ma2017lifelong} but run
on layouts that are not well-formed for them, with our dwell, charging and
parking rules added; other adaptations are possible. Finally, the console
integration displays the simulator; it is not connected to a fleet.
